\bookchapter{Глава 1: Привет, мой личный компьютер}

\chapterintro{Ваш телефон — это на самом деле самый личный суперкомпьютер в вашем кармане. Он в бесчисленное количество раз мощнее компьютеров, отправивших человека на Луну, и при этом тихо лежит в вашей ладони. Сегодня давайте как следует познакомимся с этим «молчаливым гигантом».}

\booksection{I. «Внутренние органы» телефона: глубокий разбор}

Телефон мал, но в нём есть всё. Давайте, как на уроке анатомии, поочередно познакомимся с этими точными компонентами.

\booksubsection{SoC (система на кристалле): «мозг + сердце» телефона}

SoC расшифровывается как System on Chip — «вся система на одном чипе». Это не просто CPU (центральный процессор), а интеграция CPU, GPU (графического процессора), NPU (нейронного процессора), ISP (процессора обработки изображений), модема (отвечающего за связь) и многого другого в одном кристалле.

Это как если бы в одном здании разместили мэрию, электростанцию, центр связи и полицейский участок — эффективность колоссальная.

\booksubsection{CPU: ядро универсальных вычислений}

CPU — самый центральный компонент SoC, отвечающий за обработку подавляющего большинства вычислительных задач. Мобильные CPU используют архитектуру ARM, которая отличается от архитектуры x86 настольных компьютеров. ARM-архитектура основана на сокращённом наборе команд (RISC), что даёт более низкое энергопотребление и лучше подходит для мобильных устройств.

\textbf{Дизайн CPU с большими и малыми ядрами}: современные мобильные CPU повсеместно используют архитектуру «big.LITTLE» (технология ARM), разделяя ядра на разные классы производительности:

\begin{itemize}
  \item \textbf{Сверхбольшие/большие ядра}: наивысшая производительность, отвечают за тяжёлые задачи — загрузка игр, запуск приложений, сложные вычисления. Максимальная частота (3 ГГц+), но и максимальное энергопотребление.
  \item \textbf{Средние ядра}: баланс производительности и энергопотребления, обрабатывают повседневные задачи средней нагрузки. Прокрутка ленты, просмотр видео, переключение между задачами — в основном на них.
  \item \textbf{Малые/энергоэффективные ядра}: минимальное энергопотребление, обрабатывают лёгкие фоновые задачи. Режим ожидания, музыка, приём сообщений — почти не расходуют батарею.
\end{itemize}

Ядра трёх уровней работают согласованно: планировщик системы автоматически распределяет задачи в зависимости от нагрузки — лёгкие задачи отправляются на малые ядра для экономии энергии, тяжёлые — на большие для максимальной производительности, а после выполнения снова переключаются на малые. Именно поэтому телефон может и плавно запускать игры, и оставаться в режиме ожидания целый день.

\textbf{Особенности архитектур разных брендов}:

\begin{itemize}
  \item \textbf{Qualcomm Snapdragon 8 Elite}: собственная архитектура CPU Oryon (разработана на основе приобретённой Nuvia), 2 суперъядра (4,32 ГГц) + 6 производительных ядер (3,53 ГГц), все на базе второго поколения 3-нм техпроцесса TSMC. Полный отказ от стандартных ядер ARM, скачкообразный прирост как однопоточной, так и многопоточной производительности. GPU Adreno 830 с плиточной архитектурой, прирост графической производительности на 40\%, NPU Hexagon с производительностью около 45 TOPS, прирост скорости ИИ-вывода на 45\%. (Прим.: специальная версия Samsung Galaxy имеет более высокую частоту, суперъядра до 4,47 ГГц.)
  \item \textbf{MediaTek Dimensity 9400}: 1 сверхбольшое ядро (Cortex-X925, 3,62 ГГц) + 3 производительных ядра (Cortex-X4, 3,3 ГГц) + 4 энергоэффективных ядра (Cortex-A720, 2,4 ГГц), продолжение архитектуры «все ядра большие», второй поколение 3-нм техпроцесса TSMC. GPU Immortalis-G925 (12 ядер), пиковая производительность +41\%, трассировка лучей +40\%. ИИ-процессор — NPU 890 восьмого поколения, поддержка запуска больших моделей на устройстве, LoRA-обучения на устройстве и генерации видео на устройстве.
  \item \textbf{Apple A18 Pro}: 2 производительных ядра (\textasciitilde{}4,04 ГГц) + 4 энергоэффективных ядра (\textasciitilde{}2,4 ГГц), второе поколение 3-нм техпроцесса TSMC. 6-ядерный GPU с аппаратной трассировкой лучей, 16-ядерный NPU (\textasciitilde{}35 TOPS) обеспечивает мощную вычислительную базу для ИИ на устройстве и Apple Intelligence. Собственная архитектура Apple по-прежнему лидирует по однопоточной производительности и энергоэффективности, унифицированная архитектура памяти даёт преимущество по пропускной способности для ИИ-вывода.
  \item \textbf{Huawei Kirin 9020}: установлен в серии Mate 70, собственная архитектура CPU Taishan, постоянно улучшаемые NPU и ISP, поддержка 5.5G и спутниковой связи, представляет последний прорыв китайских чипов.
\end{itemize}

\booksubsection{GPU: движок графического рендеринга}

GPU (Graphics Processing Unit) специализируется на задачах графики и параллельных вычислений. Количество его ядер намного превышает CPU (сотни и тысячи ядер), но каждое ядро функционально проще и подходит для одновременной обработки большого количества однотипных вычислений.

\textbf{Процесс графического рендеринга}:

\begin{enumerate}
  \item \textbf{Обработка вершин}: вычисление координат каждой вершины в 3D-модели.
  \item \textbf{Растеризация}: преобразование 3D-графики в пиксели на 2D-экране.
  \item \textbf{Пиксельное затенение}: вычисление цвета, текстуры и эффектов освещения для каждого пикселя.
  \item \textbf{Композиция вывода}: объединение всех слоёв в финальное изображение и отправка на дисплей.
\end{enumerate}

\textbf{Влияние GPU на повседневное использование}: не только игры! Прокрутка интерфейса, анимация, воспроизведение видео, фотофильтры, AR-приложения — всё зависит от GPU. На телефоне со слабым GPU интерфейс может терять кадры при прокрутке, а игры просто тормозят.

\textbf{GPU разных брендов}: Qualcomm использует собственный Adreno GPU; MediaTek — ARM Mali GPU; Apple — собственный GPU, который является эталоном по производительности и энергоэффективности в отрасли.

\booksubsection{NPU: движок ускорения ИИ}

NPU (Neural Processing Unit) — это процессор, специально разработанный для вычислений нейронных сетей. По сравнению с CPU и GPU, NPU обеспечивает более высокую эффективность и более низкое энергопотребление при обработке задач ИИ.

\textbf{Принцип работы нейронных вычислений}: основные операции NPU — умножение матриц и свёртка, которые являются самыми базовыми и частыми вычислениями в нейронных сетях. Благодаря аппаратной оптимизации NPU может параллельно обрабатывать большие объёмы матричных операций, работая в десятки раз быстрее CPU при значительно меньшем энергопотреблении.

\textbf{Процесс обработки ИИ}:

\begin{enumerate}
  \item \textbf{Ввод данных}: изображения с камеры, аудио с микрофона, данные с датчиков и т.д.
  \item \textbf{Предварительная обработка}: настройка формата и размера данных под требования входа нейронной сети.
  \item \textbf{Инференс (вывод)}: NPU загружает обученную модель, выполняет прямой проход по входным данным и выдаёт результат.
  \item \textbf{Постобработка}: преобразование вывода ИИ в полезные результаты — координаты лица, текст, метки сцен.
\end{enumerate}

\textbf{Практические сценарии применения}:

\begin{itemize}
  \item \textbf{Вычислительная фотография}: ночная съёмка, портретное размытие, HDR-обработка, ИИ-улучшение лица.
  \item \textbf{Голосовой помощник}: распознавание речи, понимание естественного языка, синтез речи.
  \item \textbf{Перевод в реальном времени}: ИИ-субтитры, перевод диалогов.
  \item \textbf{Интеллектуальное распознавание}: распознавание текста (OCR), распознавание объектов, распознавание сцен.
  \item \textbf{Большие модели на устройстве}: новейшие флагманские чипы уже поддерживают запуск небольших ИИ-моделей (уровня 7B параметров) на самом телефоне, обеспечивая офлайн-диалог с ИИ, реферирование документов, интеллектуальное письмо и другие функции. NPU Hexagon в Snapdragon 8 Elite имеет производительность около 45 TOPS, NPU 890 в Dimensity 9400 также поддерживает большие модели на устройстве (скорость LLM-вывода увеличена на 80\%), 16-ядерный NPU Apple A18 Pro обеспечивает около 35 TOPS.
  \item \textbf{ИИ-обработка изображений}: ИИ-удаление (удаление прохожих с фото одним нажатием), ИИ-расширение (расширение границ изображения), ИИ-перенос стиля, генерация изображений по тексту и т.д.
\end{itemize}

\booksubsection{ISP: процессор обработки изображений}

ISP обрабатывает необработанные данные с камеры, превращая световые сигналы в красивые фотографии.

\textbf{Как оценить производительность?}Смотрите на модель! Apple A18 Pro, Qualcomm Snapdragon 8 Elite, MediaTek Dimensity 9400, Huawei Kirin 9020 — это флагманский уровень. Средний уровень — Snapdragon 8s Gen3/7+ Gen3, Dimensity 8400. Не смотрите только на количество ядер — архитектура и техпроцесс (в нанометрах) важнее. 3-нм техпроцесс — самый современный, энергоэффективность значительно превосходит 5 нм и 4 нм.

\booksubsection{Оперативная память (RAM): временный рабочий стол}

Оперативная память — как ваш рабочий стол: чем больше, тем больше документов можно разложить одновременно. 12 ГБ — сегодняшний стандарт, 16 ГБ подойдёт активным пользователям, 24 ГБ и выше — выбор геймеров и энтузиастов.

Тип ОЗУ тоже важен: LPDDR5 быстрее и экономичнее LPDDR4X. У Apple ОЗУ немного, но оптимизация системы отличная, поэтому кажется, что «хватает».

\textbf{Что будет, если ОЗУ не хватает?}Система будет выгружать приложения из фона, и при переключении их придётся загружать заново. Вот почему иногда переключаешься в WeChat и обратно — а игры уже нет.

\booksubsection{Память (ROM/флеш-память): долговременный склад}

Память — как ваш книжный шкаф, место для всех фото, видео и приложений. Сейчас стандарт — UFS 3.1 или UFS 4.0, что в несколько раз быстрее старого eMMC.

\textbf{Почему важна скорость памяти?}Она определяет скорость открытия приложений, загрузки игр и сохранения фото. Со временем медленная память заставляет телефон всё больше тормозить.

\booksubsection{Экран: ваши глаза}

Экран — главное окно вашего взаимодействия с телефоном, его качество напрямую влияет на впечатления.

\begin{itemize}
  \item \textbf{Материал}: OLED — яркие цвета, высокий контраст, экономичнее (пиксели не светятся при отображении чёрного), но при низкой яркости может быть мерцание, вредное для глаз; LCD — безопаснее для глаз, дешевле, но контраст и цвета уступают OLED.
  \item \textbf{Разрешение}: 1080P — mainstream, 2K — чётче, но больше расход батареи. Концепция «Retina-дисплея» Apple означает, что с определённого расстояния глаз не различает отдельные пиксели.
  \item \textbf{Частота обновления}: 60 Гц — стандарт, 90/120 Гц — плавнее, прокрутка шёлковая. Высокая частота обновления расходует батарею, поэтому многие телефоны поддерживают адаптивную частоту (LTPO), динамически подстраиваясь под контент.
  \item \textbf{Частота опроса касаний}: влияет на отзывчивость касаний, важно для геймеров — чем выше, тем отзывчивее.
  \item \textbf{Способ регулировки яркости}: DC-регулировка безопаснее для глаз, чем PWM, при низкой яркости глаза меньше устают. Сейчас многие телефоны поддерживают высокочастотную PWM-регулировку — тоже для защиты глаз.
\end{itemize}

\booksubsection{Технология HDR экрана}

HDR (High Dynamic Range) позволяет экрану одновременно показывать более яркие светлые участки и более тёмные тени, значительно повышая детализацию изображения.

\begin{itemize}
  \item \textbf{Принцип реализации}: обычный экран имеет ограниченный диапазон яркости (SDR), разница между самым ярким и самым тёмным недостаточна. HDR-экран имеет более высокую пиковую яркость (обычно более 1000 нит), больший контраст, более широкий цветовой охват и способен отображать больше деталей.
  \item \textbf{Основные стандарты}: HDR10 (самый базовый, бесплатный и открытый), HDR10+ (продвигается Samsung, динамические метаданные), Dolby Vision (наилучшее качество, но требует лицензии). Какой стандарт поддерживает телефон, зависит от аппаратного обеспечения экрана и системной лицензии.
  \item \textbf{Реальный эффект}: при просмотре HDR-видео солнце ярче, тени глубже, цвета насыщеннее. При условии, что само видео создано в HDR и экран поддерживает HDR-отображение.
\end{itemize}

\booksubsection{Технология HDR в фотографии}

HDR при съёмке и HDR экрана — это разные вещи. Фотографический HDR решает проблему «пересветов в светлых участках и провалов в тенях» с помощью алгоритмов.

\begin{itemize}
  \item \textbf{Принцип алгоритма}: при нажатии кнопки спуска телефон фактически делает несколько снимков с разной экспозицией (недоэкспонированный сохраняет детали в светах, переэкспонированный — в тенях, нормальный — средние тона), а затем объединяет их в одно фото с высоким динамическим диапазоном.
  \item \textbf{Процесс обработки}: захват нескольких кадров компенсация выравнивания (анти-дрожание) слияние экспозиций / тональная компрессия цветокоррекция вывод финального фото.
  \item \textbf{Оптимизация эффекта}: современные телефоны уже могут делать «HDR в реальном времени» — изображение в видоискателе уже обработано HDR, что видишь, то и получаешь. Smart HDR от Apple и автоматический HDR на Android очень зрелые технологии.
\end{itemize}

\booksubsection{Аккумулятор и зарядка: источник энергии}

Аккумулятор — самая «недолговечная» деталь телефона, через два-три года заметно деградирует. Сейчас стандартная ёмкость — 5500-7500 м А·ч.

\begin{itemize}
  \item \textbf{Мощность зарядки}: в лагере Android гонка дошла до 100 Вт и даже 200 Вт — полная зарядка за десять с лишним минут; Apple всё ещё на 20-30 Вт. Вредит ли быстрая зарядка аккумулятору? С оригинальными аксессуарами проблем немного — телефон имеет термозащиту.
  \item \textbf{Беспроводная зарядка}: удобно, но низкая эффективность и сильный нагрев. Подходит для офиса — положил и заряжается, не подходит, когда спешишь.
  \item \textbf{Обратная зарядка}: телефон становится пауэрбанком для наушников и часов — на экстренный случай.
\end{itemize}

\booksubsection{Камера: магия света и тени}

Сейчас в телефонах по три-четыре камеры, у каждой своя роль:

\begin{itemize}
  \item \textbf{Основная}: самая используемая, самый большой сенсор, лучшее качество.
  \item \textbf{Сверхширокоугольная}: для пейзажей и архитектуры, более широкий обзор.
  \item \textbf{Телефото}: для съёмки удалённых объектов, бывает оптический и цифровой зум.
  \item \textbf{Макро}: для съёмки очень мелких объектов с близкого расстояния.
\end{itemize}

\textbf{Количество мегапикселей — не всё!}Размер сенсора (чем больше матрица, тем лучше), диафрагма (f/1.8 больше, чем f/2.4), стабилизация (OIS оптическая стабилизация очень важна) — всё влияет на качество. Вычислительная фотография (алгоритмы ИИ) становится всё важнее, позволяя объединять несколько снимков в один качественный.

\booksubsection{Перископический телеобъектив}

Традиционный телеобъектив требует длинного оптического пути, что делает телефон толстым. Перископический телеобъектив использует 90-градусную призму для преломления света, располагая оптический путь горизонтально внутри телефона, что позволяет достичь 5-кратного и даже 10-кратного оптического зума при сохранении тонкого корпуса.

\begin{itemize}
  \item \textbf{Принцип конструкции}: свет входит в объектив призма преломляет на 90° проходит через группу линз вдоль ширины телефона достигает сенсора. Это как положить вертикальный телескоп горизонтально.
  \item \textbf{Оптические характеристики}: высокая кратность оптического зума (5x-10x), качество намного превосходит цифровой зум. Но сенсор обычно меньше, чем у основной камеры, поэтому при слабом освещении качество среднее.
  \item \textbf{Преимущества в съёмке}: съёмка луны, удалённых зданий, крупных планов на сцене — недостижимо для обычных телефонов.
\end{itemize}

\booksubsection{OIS оптическая стабилизация vs EIS электронная стабилизация}

Стабилизация — ключевая технология, влияющая на качество фото и видео. Два метода работают по совершенно разным принципам:

\begin{itemize}
  \item \textbf{OIS (оптическая стабилизация)}: физически перемещает объектив или сенсор для компенсации дрожания рук. Внутри объектива есть микрогироскоп, определяющий направление и амплитуду дрожания, затем привод двигает группу линз в противоположном направлении. Преимущество — без потери качества, недостаток — высокая стоимость и занимает место.
  \item \textbf{EIS (электронная стабилизация)}: использует программные алгоритмы для обрезки и смещения кадра для компенсации дрожания. Телефон записывает немного большую область, а затем на основе данных гироскопа вырезает стабильную область. Преимущество — не требует дополнительного оборудования, недостаток — обрезает кадр (уменьшается поле зрения), при слабом освещении может быть размытие.
  \item \textbf{Практический выбор}: для фото в первую очередь смотрите на OIS, для видео лучше всего сочетание обоих методов. Сейчас многие флагманские телефоны поддерживают стабилизацию сдвигом сенсора «Sensor-Shift», которая эффективнее традиционной стабилизации объектива.
\end{itemize}

\booksubsection{Оптический зум vs Цифровой зум}

\begin{itemize}
  \item \textbf{Оптический зум}: физически перемещает группу линз для изменения фокусного расстояния, без потери качества. В телефонах оптический зум обычно имеет фиксированную кратность (2x, 3x, 5x, 10x) и не может плавно изменяться, как в фотоаппаратах.
  \item \textbf{Цифровой зум}: увеличивает центральную область изображения — по сути, обрезка + интерполяционное увеличение. Чем больше увеличение, тем хуже качество. 100-кратный цифровой зум выглядит круто, но реальное качество уже очень размытое.
  \item \textbf{Гибридный зум}: сочетает оптический и цифровой, используя алгоритмы многокадрового синтеза для компенсации потери качества между кратностями оптического зума. Эффект лучше чистого цифрового, но всё ещё уступает оптическому.
\end{itemize}

\booksubsection{Аудио: динамики и чип выделенной графики}

\textbf{Один динамик vs Два динамика}:

\begin{itemize}
  \item \textbf{Один динамик}: только один излучатель снизу, звук смещён в одну сторону, нет стереоэффекта. Часто встречается в бюджетных и среднебюджетных моделях.
  \item \textbf{Два динамика (стерео)}: по одному динамику сверху и снизу, создают стереозвуковое поле. При просмотре видео и прослушивании музыки левый и правый каналы чётко разделены, погружение намного сильнее. Часто разговорный динамик используется как верхний — распространённое решение.
  \item \textbf{Разница в качестве звука}: два динамика — это не просто «громче», ключевое — пространственность и многослойность. С одним динамиком при просмотре блокбастеров всегда кажется, что чего-то не хватает.
\end{itemize}

\textbf{Чип выделенной графики в игровых телефонах}:

Некоторые игровые и производительные телефоны оснащены независимым чипом обработки дисплея. Его функции:

\begin{itemize}
  \item \textbf{Интерполяция кадров}: преобразует игровой процесс с 30 fps до 60 или 120 fps с помощью алгоритмов интерполяции, делая картинку плавнее. Этот процесс выполняется на выделенном чипе, не нагружая GPU SoC.
  \item \textbf{Суперразрешение}: повышает разрешение изображения с низкого до высокого с помощью алгоритмов ИИ, делая картинку чётче.
  \item \textbf{Снижение энергопотребления}: поскольку интерполяция и суперразрешение выполняются выделенным чипом, нагрузка на SoC снижается, уменьшаются нагрев и энергопотребление, игра работает стабильнее.
\end{itemize}

\booksubsection{GPS-позиционирование: одночастотный и двухчастотный}

Позиционирование телефона в основном опирается на спутниковые навигационные системы, такие как GPS.

\begin{itemize}
  \item \textbf{Одночастотный GPS}: принимает сигнал только в диапазоне L1, точность позиционирования около 3-5 метров. Достаточно на открытой местности, но в городских каньонах с плотной высотной застройкой отражение сигналов может вызывать значительное смещение.
  \item \textbf{Двухчастотный GPS}: одновременно принимает сигналы в диапазонах L1 и L5. Диапазон L5 имеет большую длину волны и лучшую помехозащищённость, совместное решение по двум частотам устраняет большинство ионосферных ошибок, точность позиционирования достигает 1-3 метров. Даёт заметное преимущество при навигации в городе и вызове такси.
  \item \textbf{Поддержка нескольких созвездий}: современные телефоны поддерживают не только GPS, но и BeiDou (Китай), GLONASS (Россия), Galileo (Европа) и другие спутниковые системы. Чем больше спутников одновременно подключено, тем быстрее и точнее позиционирование.
\end{itemize}

\booksubsection{Датчики: маленькие щупальца телефона}

В телефоне скрыто более десятка датчиков, которые тихо работают:

\begin{itemize}
  \item \textbf{Акселерометр + гироскоп}: определяют положение телефона, поворот в играх, автоповорот экрана.
  \item \textbf{Датчик освещённости}: автоматическая регулировка яркости экрана.
  \item \textbf{Датчик приближения}: при разговоре телефон у уха — экран автоматически гаснет.
  \item \textbf{Датчик отпечатков пальцев}: сейчас mainstream — подэкранный оптический или ультразвуковой.
  \item \textbf{Распознавание лица}: Face ID от Apple — это 3D-структурированный свет, высокий уровень безопасности; многие Android — 2D-распознавание по фото, безопасность средняя.
  \item \textbf{NFC}: Near Field Communication, связь на малом расстоянии — оплата проезда, пропуск, банковская карта.
  \item \textbf{Барометр}: измеряет высоту над уровнем моря, может определить даже подъём по лестнице.
  \item \textbf{Датчик Холла}: определяет открытие/закрытие магнитного чехла.
\end{itemize}

\booksection{II. Операционные системы: Android VS iOS VS HarmonyOS}

Аппаратная часть — это тело, операционная система — это душа. У трёх систем свои особенности:

\begin{itemize}
  \item \textbf{Android}: как открытое сообщество — производители телефонов могут глубоко кастомизировать систему, выбор огромен. Можно менять темы, устанавливать сторонние магазины приложений, экспериментировать. Но из-за открытости — сильная фрагментация, опыт на разных брендах сильно различается.
  \item \textbf{iOS}: как тщательно ухоженный сад — Apple контролирует и софт, и железо, опыт унифицированный и плавный. Качество приложений высокое, защита приватности хорошая, но свобода ограничена — многое нужно делать по правилам Apple.
  \item \textbf{HarmonyOS}: новый вид от Huawei, ключевая особенность — распределённые возможности, позволяющие разным устройствам бесшовно взаимодействовать. Подробнее расскажем позже.
\end{itemize}

Экосистемы всех трёх чрезвычайно богаты — выбирайте тот образ жизни, который вам по душе. Нет абсолютно хорошего или плохого, есть только подходящее или нет.

\booksection{III. Полный разбор интерфейсов телефона}

Интерфейсы телефона — это мосты, соединяющие его с внешним миром. Разные интерфейсы сильно отличаются:

\booksubsection{Micro USB}

\begin{itemize}
  \item \textbf{Физические характеристики}: плоский трапециевидный разъём, различает верх и низ (неправильно не вставить). Практически вышел из употребления в телефонах, но всё ещё используется в некоторых бюджетных и старых устройствах.
  \item \textbf{Ограничения передачи}: поддерживает только стандарт USB 2.0, максимальная скорость передачи 480 Мбит/с, мощность зарядки обычно не превышает 10 Вт.
  \item \textbf{Недостатки}: несимметричный разъём, малый ресурс подключений, низкая скорость передачи, малая мощность зарядки.
\end{itemize}

\booksubsection{Lightning (фирменный разъём Apple, снят с производства)}

\begin{itemize}
  \item \textbf{Физические характеристики}: 8-контактный разъём собственной разработки Apple, симметричный — можно вставлять любой стороной. Использовался в iPhone 14 и более ранних моделях, полностью выведен из употребления.
  \item \textbf{Ограничения передачи}: физически разъём Lightning поддерживает только стандарт USB 2.0, максимальная скорость 480 Мбит/с. Даже при использовании кабеля USB-C на Lightning скорость передачи ограничена самим разъёмом Lightning.
  \item \textbf{Причина вывода из употребления}: ЕС принял закон, требующий с 2024 года использовать единый разъём USB-C для всех телефонов. Apple начиная с серии iPhone 15 полностью перешла на Type-C. Разъём Lightning официально уходит в историю.
  \item \textbf{Примечание}: кабели с сертификацией Apple MFi содержат чип аутентификации, несертифицированные кабели могут не заряжать или не передавать данные. Если вы всё ещё пользуетесь старым iPhone, рекомендуется иметь запасной кабель Lightning.
\end{itemize}

\booksubsection{USB Type-C}

\begin{itemize}
  \item \textbf{Физические характеристики}: овальный 24-контактный разъём, симметричный. Стал стандартным разъёмом для телефонов и большинства электронных устройств. ЕС законодательно обязал с 2024 года использовать USB-C для всех телефонов, планшетов, наушников и других портативных устройств, другие регионы мира также следуют этому примеру.
  \item \textbf{Протокол передачи}: Type-C — это только форма разъёма, реальная скорость передачи зависит от внутреннего протокола. Поддерживает USB 2.0 (480 Мбит/с), USB 3.0/3.1/3.2 (5-20 Гбит/с) и даже USB4 (40 Гбит/с).
  \item \textbf{Важное различие}: не все разъёмы Type-C поддерживают USB 3.0! Многие телефоны используют разъём Type-C, но внутри подключены только контакты USB 2.0, и скорость передачи такая же медленная, как у старых разъёмов. Стандартные версии iPhone 16/17 поддерживают только скорость USB 2.0, версии Pro — USB 3.0 (требуется кабель USB-C на USB-C).
  \item \textbf{Возможности зарядки}: поддержка протоколов быстрой зарядки, таких как USB PD, максимальная мощность до 240 Вт.
  \item \textbf{Видеовыход}: поддержка DisplayPort Alt Mode, можно выводить видео на монитор через переходник.
\end{itemize}

\booksubsection{3,5-мм аудиоразъём}

\begin{itemize}
  \item \textbf{Физические характеристики}: круглый 3- или 4-сегментный штекер, 3 сегмента передают левый и правый аудиоканалы, 4 сегмента добавляют канал микрофона.
  \item \textbf{Особенности}: передача аналогового сигнала, подключил и работает, не требует драйверов или преобразования. Стабильное качество звука, нулевая задержка.
  \item \textbf{Текущее состояние}: большинство флагманских телефонов отказались от 3,5-мм разъёма, требуется переходник Type-C или использование Bluetooth-наушников. Некоторые среднебюджетные и игровые телефоны всё ещё сохраняют его.
\end{itemize}

\booksection{IV. Система продуктов телефонов и правила именования}

Глядя на ослепительное разнообразие продуктовых линеек разных брендов, понимание правил именования поможет быстро определить подходящую модель.

\booksubsection{Продуктовые линейки основных брендов:}

\begin{itemize}
  \item \textbf{Apple}: iPhone (стандартная версия) / iPhone Plus (версия с большим экраном) / iPhone Pro (профессиональная версия) / iPhone Pro Max (профессиональная версия с большим экраном) / iPhone SE (бюджетная версия) / iPhone 16e (бюджетная ИИ-версия с чипом A18 и собственным модемом Apple C1). В 2026 году серия iPhone 17 полностью оснащена чипом A19 Pro, возможности ИИ ещё больше расширены.
  \item \textbf{Huawei}: серия Mate (бизнес-флагман) / серия P/Pura (фотофлагман) / серия Mate X (складной флагман) / серия nova (молодёжная среднебюджетная) / серия Enjoy (начальный уровень)
  \item \textbf{Xiaomi}: цифровая серия Xiaomi (флагман) / Xiaomi MIX (исследовательский флагман/складной) / серия Redmi K (производительный флагман) / серия Redmi Note (средний уровень) / цифровая серия Redmi (начальный уровень)
  \item \textbf{OPPO}: серия Find X (флагман) / серия Reno (средний уровень) / серия A (начальный уровень) / цифровая серия OnePlus (производительный флагман)
  \item \textbf{vivo}: серия X (фотофлагман) / цифровая серия iQOO (производительный флагман) / серия S (средний уровень с акцентом на селфи) / серия Y (начальный уровень)
  \item \textbf{Samsung}: серия Galaxy S (флагман) / Galaxy Z Fold/Flip (складной флагман) / серия Galaxy A (средний уровень)
\end{itemize}

\booksubsection{Значение распространённых суффиксов:}

\begin{itemize}
  \item \textbf{Pro}: сокращение от Professional, профессиональная версия. Обычно обновлены камера, экран, процессор и другие ключевые компоненты по сравнению со стандартной версией.
  \item \textbf{Max}: максимальная версия. Самый большой экран и аккумулятор, но конфигурация не обязательно выше, чем у Pro.
  \item \textbf{Ultra}: ультра-версия. Обычно представляет максимальную конфигурацию в серии, максимальная начинка и самая высокая цена. Например, Samsung S24 Ultra, Xiaomi 14 Ultra.
  \item \textbf{SE}: Special Edition, специальное издание, на практике — бюджетная версия. Урезанная конфигурация, сниженная цена. Например, iPhone SE.
  \item \textbf{Plus}: увеличенная версия. Экран больше стандартного, аккумулятор тоже, но конфигурация в основном та же.
  \item \textbf{Lite / Youth Edition}: облегчённая версия. Упрощённая конфигурация, более низкая цена.
  \item \textbf{Note}: изначально обозначало телефон с большим экраном, сейчас разные бренды используют по-разному. Redmi Note — это среднебюджетная серия.
  \item \textbf{Neo}: новая версия. Обычно обозначает среднебюджетную модель с акцентом на производительность.
\end{itemize}

\booksection{V. Модем и антенна: невидимые герои связи телефона}

Телефон может звонить и выходить в интернет благодаря тандему модемного чипа и антенны.

\booksubsection{Модемный чип:}

Модем — это процессор в телефоне, отвечающий за беспроводную связь. Он модулирует цифровые сигналы в радиоволны для передачи и демодулирует принятые радиоволны обратно в цифровые сигналы. Производительность модема напрямую влияет на качество звонков, скорость интернета и стабильность сигнала. Qualcomm — крупнейший поставщик модемов для телефонов, MediaTek и Samsung также имеют собственные модемы. Apple, начиная с iPhone 16e, использует собственный модемный чип C1, что знаменует официальный вход Apple в эру собственных модемов; чип C1 выполнен по 4-нм техпроцессу с отличным энергопотреблением. iPhone 16/16 Plus/16 Pro/16 Pro Max всё ещё используют модем Qualcomm (Snapdragon X75 5G), ожидается, что последующие модели будут постепенно переходить на собственные модемы. Более ранние серии iPhone XS/11 использовали модемы Intel с плохим приёмом сигнала, что было проблемой для многих пользователей.

\booksubsection{Дизайн антенны:}

Антенны телефона спрятаны внутри корпуса и обычно расположены по периметру рамки. Дизайн антенны напрямую влияет на качество приёма сигнала:

\begin{itemize}
  \item \textbf{Технология MIMO}: Multiple-Input Multiple-Output, использует несколько антенн для одновременной передачи и приёма данных, повышая скорость и стабильность. 4×4 MIMO быстрее, чем 2×2 MIMO.
  \item \textbf{Количество антенн}: флагманские модели обычно имеют более десятка антенн (сотовая связь, Wi-Fi, Bluetooth, NFC, GPS — каждой нужна своя антенна).
  \item \textbf{Влияние металлического корпуса}: металл экранирует сигнал, поэтому телефоны с металлическим корпусом требуют антенных полос (пластиковые вставки на рамке). Стеклянные и пластиковые корпуса меньше влияют на сигнал.
\end{itemize}

\booksubsection{Эволюция сетевых технологий:}

\begin{itemize}
  \item \textbf{4G LTE}: теоретическая пиковая скорость — нисходящий канал 1 Гбит/с / восходящий 500 Мбит/с (Cat.18), реальный опыт обычно 50-150 Мбит/с. Самое широкое покрытие, достаточно.
  \item \textbf{5G NR}: теоретическая пиковая скорость — нисходящий канал 10 Гбит/с+. Диапазон Sub-6GHz имеет хорошее покрытие, mmWave (миллиметровые волны) — экстремальная скорость, но очень малая зона покрытия.
  \item \textbf{5G-A (5G Advanced)}: усиленная версия 5G, также называется 5.5G. Скорость нисходящего канала до 10 Гбит/с, восходящего — 1 Гбит/с, плотность подключений и точность позиционирования значительно улучшены. С 2024 года China Mobile, China Telecom и China Unicom начали развёртывание сетей 5G-A в ключевых городах. Флагманские телефоны, такие как серии Huawei Mate 70/80 и Xiaomi 15/16, уже поддерживают 5G-A. 5G-A также приносит возможность «интегрированного восприятия и связи» — базовая станция может одновременно осуществлять связь и воспринимать окружающую среду, обеспечивая инфраструктуру для низковысотной экономики (доставка дронами) и V2X.
\end{itemize}

\booksubsection{Классификация сетевых диапазонов:}

Поддерживаемые телефоном сетевые диапазоны определяют, в каких странах и у каких операторов он может работать:

\begin{itemize}
  \item \textbf{FDD-LTE Band (серия B)}: например B1, B3, B5 и т.д. — диапазоны 4G. Разные страны используют разные диапазоны, Китай в основном использует B1/B3/B5/B8/B34/B38/B39/B40/B41.
  \item \textbf{TDD-LTE Band}: другой режим 4G, Китай широко использует B38/B39/B40/B41.
  \item \textbf{5G NR Band (серия N)}: например n1, n3, n28, n41, n78, n79 и т.д. n78 (3,5 ГГц) — основной диапазон 5G в Китае, n41 (2,5 ГГц) больше используется China Mobile.
  \item \textbf{Совет при покупке}: если планируете использовать телефон за границей, перед покупкой проверьте, поддерживает ли телефон диапазоны целевой страны. Китайские версии телефонов обычно имеют наиболее полную поддержку китайских диапазонов, но могут не иметь некоторых зарубежных диапазонов.
\end{itemize}

\booksection{VI. Сколько ролей на самом деле играет ваш телефон?}

Он давно уже не просто телефон:

\begin{enumerate}
  \item \textbf{Средство связи}: звонки, SMS, WeChat, электронная почта…
  \item \textbf{Умная камера}: фото, видео, сканирование документов…
  \item \textbf{Кошелёк}: оплата, проездной, банковские карты, карты лояльности…
  \item \textbf{Информационный портал}: новости, погода, акции, знания…
  \item \textbf{Игровая консоль}: от головоломок до AAA-игр…
  \item \textbf{Мобильный офис}: документы, таблицы, совещания, совместная работа…
  \item \textbf{Помощник здоровья}: шагомер, пульс, сон, спорт…
  \item \textbf{Пульт умного дома}: управление светом, кондиционером, телевизором, роботом-пылесосом…
  \item \textbf{Цифровое удостоверение личности}: электронный ID, электронные права, код здоровья…
  \item \textbf{Развлекательный центр}: музыка, видео, электронные книги, подкасты…
\end{enumerate}

Подумайте только — сколько всего умещается в одном маленьком устройстве! Поэтому научиться хорошо им пользоваться — значит научиться жить лучше.

\begin{tcbpractice}{Практика: узнайте конфигурацию оборудования}
В [Настройках] посмотрите, можете ли вы найти модель процессора, объём оперативной памяти (RAM) и встроенной памяти (ROM). Составьте представление о возможностях вашего телефона. Пользователи Android могут скачать «AnTuTu» или «DevCheck» для просмотра подробной информации об оборудовании.

\end{tcbpractice}
